% !PRINTED 162
\weberpage{162}
%<<CONT
giebt unter dieser Voraussetzung für die gemeinsame Wurzel
einen rationalen Ausdruck durch die Coëfficienten.

Setzt man diesen Ausdruck in $\psi(x)$ ein, so erhält man,
wenn man alle Nenner fortschafft, eine ganze rationale Function
der Coëfficienten $a$, $b$, $c$, die wir mit $S$ bezeichnen wollen, und
% [corrected 2026-08-07, 10x zoom strip 2/10] The Sperrung starts mid-paragraph at "deren
% Verschwinden" and runs to the end of the sentence -- an unnumbered result spaced inside
% running prose, the class the spec previously denied existed. Raised at high confidence by
% this page's translation pass and confirmed.
\gesperrt{deren Verschwinden in Verbindung mit $R = 0$ die nothwendige und
hinreichende Bedingung dafür enthält, dass die drei Gleichungen (1) eine gemeinsame Wurzel
haben.}

Ersetzen wir nun die Coëfficienten $a$, $b$, $c$ der Functionen
(1) durch ganze rationale Functionen einer zweiten Unbekannten
$y$, deren Coëfficienten mit $\alpha$, $\beta$, $\gamma$ bezeichnet und vorläufig als
unabhängige Veränderliche betrachtet werden mögen.

$R$ und $S$ werden dann ganze rationale Functionen von
$\alpha$, $\beta$, $\gamma$ und $y$ und wenn wir also die Gleichungen

\weq{2}{R = 0, \qquad S = 0}

als Gleichungen für $y$ betrachten, so können wir nach den
Regeln des vorigen Paragraphen ihre Resultante $T$ bilden, die,
wenn sie nicht identisch verschwindet, eine ganze rationale
Function der Coëfficienten $\alpha$, $\beta$, $\gamma$ sein wird, und die Gleichung

\weq{3}{T = 0}

ist die nothwendige und hinreichende Bedingung dafür, dass das
Gleichungssystem (1) durch ein oder mehrere Werthepaare $x$, $y$
% [corrected 2026-08-07, 10x zoom strips 6/10 and 7/10] Two spaced runs here, both missed.
% "Resultante des Systems (1)." carries the Sperrung through the reference number. The second
% run starts at "Elimination" -- "Man kann also aus den drei Gleichungen (1) durch" is set
% normally, and the break falls at "durch | Elimi-". The translation pass called both, and
% called the second one's starting word correctly.
befriedigt werden kann. $T$ heisst die \gesperrt{Resultante des Systems
(1).} % [source-original logical slip, confirmed 2026-08-25 against scan p. 187 / printed p. 162]
% The print says that eliminating x and y forms a final equation "für x", although x has just
% been eliminated; the next paragraph introduces z and says T becomes an equation for z.
% Preserved verbatim rather than silently substituting a conjectural variable.
Man kann also aus den drei Gleichungen (1) durch \gesperrt{Elimination der Unbekannten $x$ und $y$ eine Endgleichung für
$x$ bilden.}

Ersetzt man nun die $\alpha$, $\beta$, $\gamma$ wieder durch ganze rationale
Functionen einer dritten Unbekannten $z$, so geht $T$ in eine
Gleichung für $z$ über. Die Functionen $f$, $\varphi$, $\psi$ werden Functionen
der drei Unbekannten $x$, $y$, $z$. Der Grad der Gleichung (3) giebt
die Anzahl der Werthsysteme $x$, $y$, $z$, für die die drei Gleichungen
(1) zugleich befriedigt sind.

Sind aber die Coëfficienten so beschaffen, dass $T$ identisch
(für alle Werthe von $z$) verschwindet, so giebt es unendlich viele
Werthsysteme $x$, $y$, $z$, die den Gleichungen (1) zugleich genügen.

Betrachtet man $x$, $y$, $z$ als Cartesische Coordinaten, so
stellen die Gleichungen (1), jede für sich, eine Oberfläche dar.
Der Grad der Endgleichung $T = 0$ in Bezug auf $z$ giebt die
%CONT>>
